\section{Conclusion, Development Context, and Reflection}
\label{sec:conclusion}

This project reproduced the core idea of GuP through a correctness-first implementation for exact subgraph matching and organized the report to match the required grading items explicitly. The final report now contains a literature taxonomy, a clear reproduction-scope statement, mandatory experiments on three successful datasets, and a filtering-strategy ablation study.

From the perspective of the development history of subgraph matching, the field has evolved from earlier candidate-space and access-method ideas~\cite{he2008graphql} toward increasingly integrated systems that combine filtering, order optimization, conflict reuse, and execution engineering~\cite{han2019daf,sun2020study,sun2021rapidmatch}. GuP is an instructive modern point in this development line because it emphasizes dynamic, state-aware pruning while remaining conceptually decomposable into interpretable modules.

Our experimental results lead to three final conclusions. First, the project now satisfies the mandatory multi-dataset requirement with successful experiments on \texttt{Yeast}, \texttt{Human}, and \texttt{HPRD}. Second, reservation-based pruning is the only reproduced GuP mechanism that contributes clear and repeated search-space reductions on the current workloads. Third, the current Python prototype still has a visible scalability boundary, which is why \texttt{WordNet} and \texttt{Patents} remain time-boxed boundary datasets rather than full successful result datasets.

The most important lesson for us is that reproducing a research paper is not only about reproducing the same final metric values. It also requires understanding which algorithmic ideas are central, which engineering details control practical performance, and how to report scope boundaries honestly. In this project, the largest gap between the paper and the prototype lies not in correctness, but in the strength of the full guarded candidate-space machinery and the efficiency of the nogood mechanism. This gave us a concrete understanding of the difference between an algorithmic idea that is valid in principle and a system implementation that is competitive in practice.

Overall, we regard the project as a successful course-project-level reproduction of the core GuP method. It demonstrates that guard-based pruning, especially reservation-based pruning, can reduce the search space on real workloads, while also making clear what remains to be improved in a more faithful future reproduction.
